\section{Numerical Integration}

\subsection{Newton-Cotes Integration}

\begin{exercise}\label{ex:integrals}
	For each integral below, use the trapezoidal rule to estimate the integral
	for $m=10\cdot 2^k$ nodes for $k=1,2,\ldots,10$. Make a log-log plot of the
	errors and confirm or refute second-order accuracy.
	\begin{enumerate}
		\item $\displaystyle \int_0^1 x\log(1+x)\, dx = \frac{1}{4}$

		\item $\displaystyle \int_0^1 x^2 \tan^{-1}x\, dx = \frac{\pi-2+2\log 2}{12}$

		\item $\displaystyle \int_0^{\pi/2}e^x \cos x\, dx = \frac{e^{\pi/2}-1}{2}$

		\item $\displaystyle \int_0^1 {\tan^{-1}(\sqrt{2+x^2})\over
				      (1+x^2)\sqrt{2+x^2}}dx={5\pi^2\over 96}$

		\item $\displaystyle \int_0^1 \sqrt{x} \log(x) \, dx = -\frac{4}{9}$
		      (Hint: Although the integrand has the limiting value zero as $x\to 0$,
		      it cannot be evaluated naively at $x=0$. One possibility is to start the
		      integral at $x=\epsilon_{\rm mach}$. Alternatively, you can define a
		      wrapper function $g(x)$ that returns $g(0)=0$ and $g(x)=\sqrt{x}\log(x)$
		      otherwise.)

		\item $\displaystyle \int_0^1 \sqrt{1-x^2}\,\, dx = \frac{\pi}{4}$
	\end{enumerate}
\end{exercise}

\begin{solution}
	To evaluate the given definite integrals, we employ the composite
	trapezoidal rule. For an interval $[a, b]$ discretized into $m$ equally
	spaced sub-intervals with step size $h = \frac{b-a}{m}$, the composite
	trapezoidal rule approximations are given by:
	\[
		\int_a^b f(x) \, dx \approx \frac{h}{2} \left( f(x_0) +
		2\sum_{i=1}^{m-1} f(x_i) + f(x_m) \right)
	\]
	where $x_i = a + i h$. The number of nodes was systematically increased
	according to $m = 10 \cdot 2^k$ for $k = 1, 2, \ldots, 10$, yielding a
	maximum of $10,240$ nodes.

	\paragraph{Results}
	Table~\ref{tab:integral_results} summarizes the final numerical
	approximations achieved at $k=10$ ($10,240$ nodes) alongside their
	theoretical exact values.

	\begin{center}
		\centering
		\captionof{table}{Comparison of exact analytical solutions and numerical
			approximations evaluated at $m = 10,240$ nodes.}
		\label{tab:integral_results}
		\begin{tabular}{clll}
			\toprule
			\textbf{Integral} & \textbf{Exact Value}   & \textbf{Numerical Approximation} & \textbf{Absolute Error} \\
			\midrule
			1                 & $0.25$                 & $0.2500000009482$                & $9.48 \times 10^{-10}$  \\
			2                 & $\approx 0.210657251$  & $0.2106572528715$                & $1.87 \times 10^{-9}$   \\
			3                 & $\approx 1.905238690$  & $1.9052386790888$                & $1.14 \times 10^{-8}$   \\
			4                 & $\approx 0.514041897$  & $0.5140418956028$                & $1.86 \times 10^{-9}$   \\
			5                 & $\approx -0.444444444$ & $-0.4444422428639$               & $2.20 \times 10^{-6}$   \\
			6                 & $\approx 0.785398163$  & $0.7853978796777$                & $2.84 \times 10^{-7}$   \\
			\bottomrule
		\end{tabular}
	\end{center}

	\begin{center}
		\centering
		\includegraphics[width=0.9\linewidth]{figures/5-1/newton-cotes-results.pdf}
		\captionof{figure}{Evaluation of the integrals as a function of the
			number of nodes $m$. The exact analytical solutions are denoted by
			dashed horizontal lines.}
		\label{fig:results}
	\end{center}

	\begin{center}
		\centering
		\includegraphics[width=0.9\linewidth]{figures/5-1/newton-cotes-absolute-error.pdf}
		\captionof{figure}{Log-log progression of the absolute error relative to
			the number of nodes $m$. Note the disparate convergence behaviors
			between smoothly behaving integrands (1-4) and those with singular
			derivatives (5-6).}
		\label{fig:abs_error}
	\end{center}

	\paragraph{Discussion \& Convergence Analysis}
	The error bound for the composite trapezoidal rule is theoretically given by:
	\[
		E \leq \frac{(b-a)}{12} h^2 \max_{x \in [a,b]} |f''(x)| = \mathcal{O}(h^2)
	\]
	This stipulates that the algorithm demonstrates second-order accuracy
	\textit{only} if the integrand $f(x)$ is twice continuously differentiable,
	$f \in C^2[a,b]$.

	As observed in Table~\ref{tab:integral_results} and
	Figure~\ref{fig:abs_error}, all the integrals align with theoretical
	expectations, exhibiting strict second-order accuracy. In a log-log plot of
	error versus nodes $m$, this manifests as a line with a slope of $-2$.

\end{solution}

\begin{exercise}
	For each integral in Exercise~\ref{ex:integrals} above, apply the Simpson formula and compare the errors to fourth-order convergence.
\end{exercise}

\subsection{Free-nodes Integration}

\begin{exercise}
	Write a program that computes the nodes and weights for the Gauss-Legendre quadrature in the interval $[-1,1]$. Below are some tips for the implementation.

	\begin{itemize}
		\item To find the roots $x_k$, with $k=1,\ldots,n$ of the Legendre polynomial $P_n(x)$ of order $n$ use Newton's method. Take as the initial condition for the $k$th root,
		      \[
			      x^{(0)}_k=\cos(\phi_k)\,,
			      \qquad
			      \phi_k={4k-1\over 4n+2}\pi\,.
		      \]
		\item In order to use Newton's method you need to evaluate $P_n(x_k^{(i)})$ and the derivative $P'_n(x_k^{(i)})$, where $x_k^{(i)}$ is the $i$th iteration in the search for the $k$th root $x_k$. You can obtain these by first solving the recursion relation,

		      \[
			      (m+1)P_{m+1}(x) = (2m+1)xP_{m}(x) - mP_{m-1}(x)\,,
		      \]

		      for $m=1,\ldots,n-1$, where $P_0(x) = 1$ and $P_1(x) = x$, with $x=x^{(i)}_k$. The derivative is then computed using the relation,

		      \[
			      (x^{2}-1)P'_{n}(x)=n[xP_{n}(x)-P_{n-1}(x)]\,.
		      \]

		\item Once you reach the root $x_k$ within some tolerance, the value of the derivative $P'_n(x_k)$ allows you to obtain the corresponding weight using

		      \[
			      w_k= {2\over (1-x_k^2)[P'_n(x_k)]^2}\,.
		      \]
	\end{itemize}

\end{exercise}

\subsection{Advantage topics in integration}

\begin{exercise}
	For each integral, compute approximations using Gauss-Legendre quadrature rules with $n=4,6,8,\ldots,40$. Plot the error $|I_n-I|$ as a function of $n$ on a semi-log scale.
	\begin{enumerate}
		\item $\int_{-1}^1 e^{-4x}\, dx = \sinh(4)/2$
		\item $\int_{-1}^1 e^{-9x^2} = \sqrt{\pi}\, \operatorname{erf}(3)/3$
		\item $\int_{-1}^1 \operatorname{sech}(x) \, dx = 2 \tan^{-1}[ \sinh (1) ]$
		\item $\int_{-1}^1 \frac{1}{1+9x^2}\, dx = \frac{2}{3} \tan^{-1}(3)$
		\item $\int_{\pi/2}^{\pi} x^2 \sin 8x \, d x = -\frac{3 \pi^2}{32}$
	\end{enumerate}
\end{exercise}

\begin{exercise}
	For each improper integral, compute approximations using the proper double-exponential quadrature rule with $N=n/2$ and $n=4,6,8,\ldots,60$. Plot the errors as functions of $n$ on a semi-log scale.
	\begin{enumerate}
		\item $\int_{-\infty}^\infty \dfrac{1}{1+x^2+x^4}\, dx = \dfrac{\pi}{\sqrt{3}}$
		\item $\int_{-\infty}^\infty e^{-x^2}\cos(x)\, dx = e^{-1/4}\sqrt{\pi}$
		\item $\int_{-\infty}^\infty (1+x^2)^{-2/3}\, dx = \dfrac{\sqrt{\pi}\,\Gamma(1/6)}{\Gamma(2/3)}$
		\item $\int_{0}^{\infty} {1\over 1+x^2}\,dx = \frac{\pi}{2}$
		\item $\int_0^\infty \frac{e^{-x}}{\sqrt{x}}\,dx = \sqrt{\pi}$
	\end{enumerate}
\end{exercise}

\begin{exercise}
	For each integral, compute approximations using the Gauss-Legendre quadrature rule with $n=4,6,8,\ldots,60$, and the proper double-exponential quadrature rule with $N=n/2$. Plot the errors as functions of $n$ on a semi-log scale.
	\begin{enumerate}
		\item $\int_{0}^1 \frac{1}{\sqrt{1-x^2}}\, dx = \dfrac{\pi}{2}$
		\item $\int_0^1 \sqrt{x} \log(x) \, dx = -\frac{4}{9}$
		\item $\int_0^1 \sqrt{1-x^2}\,\, dx = \frac{\pi}{4}$
		\item $\int_{0}^1 (\log x)^2\, dx = 2$
		\item $\int_{0}^{\pi/2} \log(\cos(x))\, dx = -\frac{\pi}{2}\log(2)$
		\item $\int_{0}^{\pi/2} \sqrt{\tan(x)}\, dx = \dfrac{\pi}{\sqrt{2}}$
	\end{enumerate}
\end{exercise}
